\documentclass[11pt,a4paper]{article}
\usepackage[utf8]{inputenc}
\usepackage{amsmath,amssymb,amsthm}
\usepackage[margin=1in]{geometry}
\usepackage{hyperref}

\newtheorem{theorem}{Theorem}

\title{\textbf{Derivation D6: Statistical Confidence Bounds and Finite-Size Scaling Extrapolation}}
\author{\textbf{Bhogaraju M Pranav Rohit} (Roll No: 250001014) \\ Analyst --- Team Avengers (Project 06)}
\date{November 2026}

\begin{document}
\maketitle

\section{Statement}
\textit{Derivation of 95\% confidence intervals ($\pm 1.96 \frac{s}{\sqrt{T}}$) and finite-size scaling extrapolation ($p_c(L) = p_c + a \cdot L^{-0.75}$).}

\section{Part 1: Derivation of the 95\% Confidence Interval}

Let $X_1, X_2, \dots, X_T$ denote the percolation thresholds measured across $T$ independent Newman--Ziff simulation sweeps on an $L \times L$ lattice using distinct pseudo-random seeds.
\begin{itemize}
    \item The samples $X_i$ are independent and identically distributed with population mean $\mu = \mathbb{E}[X]$ and finite variance $\sigma^2 = \text{Var}(X)$.
    \item The sample mean and sample variance are:
    \[
    \bar{X}_T = \frac{1}{T}\sum_{i=1}^T X_i, \quad s^2 = \frac{1}{T-1}\sum_{i=1}^T (X_i - \bar{X}_T)^2
    \]
\end{itemize}

By the Central Limit Theorem, as $T \ge 30$, the standardized distribution converges:
\[
Z = \frac{\bar{X}_T - \mu}{s / \sqrt{T}} \xrightarrow{d} \mathcal{N}(0, 1)
\]
For a two-sided $95\%$ confidence level, the significance parameter is $\alpha = 0.05$. We find the critical normal value $z_{1 - \alpha/2} = z_{0.975}$ such that:
\[
\Phi(z_{0.975}) - \Phi(-z_{0.975}) = 0.95 \implies z_{0.975} \approx 1.95996 \approx 1.96
\]
Substituting $Z$:
\[
\mathbb{P}\left(-1.96 \le \frac{\bar{X}_T - \mu}{s / \sqrt{T}} \le 1.96\right) = 0.95
\]
Rearranging terms around the true mean $\mu$:
\[
\mathbb{P}\left(\bar{X}_T - 1.96 \frac{s}{\sqrt{T}} \le \mu \le \bar{X}_T + 1.96 \frac{s}{\sqrt{T}}\right) = 0.95
\]
The confidence interval half-width is $\Delta = 1.96 \frac{s}{\sqrt{T}}$.

\section{Part 2: Finite-Size Scaling Extrapolation}

Continuous phase transitions in statistical mechanics follow finite-size scaling theory.
Near the infinite-lattice percolation threshold $p_c(\infty)$, the correlation length $\xi$ diverges as a power law:
\[
\xi \propto |p - p_c|^{-\nu}
\]
where $\nu$ is the correlation-length universality exponent. For 2D square-lattice percolation, $\nu = 4/3$.

On a finite lattice of linear dimension $L$, the maximum correlation length is truncated by the lattice boundary: $\xi \sim L$.
Equating $|p_c(L) - p_c(\infty)|^{-\nu} \sim L$ and taking the inverse power:
\[
p_c(L) - p_c(\infty) \propto L^{-1/\nu}
\]
Substituting the 2D universality exponent $\nu = 4/3$:
\[
-\frac{1}{\nu} = -\frac{1}{4/3} = -\frac{3}{4} = -0.75
\]
\[
p_c(L) = p_c(\infty) + a \cdot L^{-0.75}
\]
where $a$ is an amplitude constant.
Plotting empirical means $\bar{p}_c(L)$ against the transformed coordinate $x = L^{-0.75}$ yields a linear relationship. In the thermodynamic limit $L \to \infty$, $x \to 0$. Linear regression to the y-intercept extracts $p_c(\infty) \approx 0.592746$.

\section{Code Mapping}
\begin{itemize}
    \item \textbf{C++ Scaling Driver:} \texttt{app/extrapolate.cpp} computes $\bar{p}_c(L)$, $s$, and $\pm 1.96 \frac{s}{\sqrt{T}}$.
    \item \textbf{Figure Generation Pipeline:} \texttt{bench/plot\_results.py} computes the linear regression of $\bar{p}_c(L)$ against $L^{-0.75}$ to extrapolate $p_c(\infty)$.
\end{itemize}

\section{References}
\begin{enumerate}
    \item Stauffer, D., \& Aharony, A. (2018). \textit{Introduction to Percolation Theory} (2nd ed., Chapter 4: Finite-Size Scaling). Taylor \& Francis.
    \item Newman, M. E. J., \& Ziff, R. M. (2000). Efficient Monte Carlo algorithm and high-precision results for percolation. \textit{Physical Review Letters}, 85(19), 4104--4107.
\end{enumerate}

\end{document}
